\section{General geometry and observable transport}
\label{sec:general-geometry}
The Meyer audit answers a question about one operator. It does not identify
what makes transport discoverable in arbitrary geometry. We now separate
three objects: the geometry that admits an iteration, the observables in
which that iteration has a short recurrence, and the cost of acquiring that
recurrence. Curvature is not itself the obstruction. Finite observable closure
is a classical Koopman concept~\cite{brunton,ssd}; our purpose here is to state
its mirror-compatibility restrictions and the discovery contract explicitly.
The following results concern smooth interior mirror descent with fixed step
size. Nonsmooth proximal maps and retained-memory algorithms still fit the
complete-state observable formulation, but not the smooth gradient test below.

\subsection{Which dynamics can a specified mirror geometry admit?}
Let $\Omega\subset\R^n$ be open and simply connected, let $h\in C^2(\Omega)$
have positive-definite Hessian, and assume $\nabla h$ is injective. Let
$F:\Omega\to\Omega$ be $C^1$ and $\eta>0$. A mirror step for a differentiable
objective satisfies
\begin{equation}
 \nabla h(F(x))=\nabla h(x)-\eta\nabla f(x).
 \label{eq:general-mirror}
\end{equation}
\begin{theorem}[Geometry--dynamics compatibility]
There exists a $C^2$ objective $f$ satisfying \eqref{eq:general-mirror} if and
only if $H_h(F(x))DF(x)$ is symmetric everywhere. If $\Omega$ is also convex,
that objective is convex if and only if
\begin{equation}
 H_h(x)-H_h(F(x))DF(x)\succeq0 \quad(x\in\Omega).
 \label{eq:general-convexity}
\end{equation}
The objective is unique up to an additive constant on a connected domain.
\end{theorem}
\begin{proof}
Define $g(x)=[\nabla h(x)-\nabla h(F(x))]/\eta$. Differentiation gives
\[ Dg(x)=\eta^{-1}[H_h(x)-H_h(F(x))DF(x)]. \]
Since $H_h(x)$ is symmetric,
$Dg$ is symmetric precisely under the stated condition. The closed one-form
$g\cdot dx$ is exact on a simply connected domain, giving $g=\nabla f$.
The converse follows by differentiating \eqref{eq:general-mirror}.
On a convex domain a $C^2$ function is convex precisely when its Hessian is
positive semidefinite. Injectivity of $\nabla h$ identifies the prescribed
$F$ with the mirror update wherever its image lies in $\Omega$.
\end{proof}
This also tests a candidate finite representation: if $\psi$ is an embedding
with decoder $\chi$ and $T\psi(\Omega)\subseteq\psi(\Omega)$, set
$F=\chi T\psi$ and apply the theorem. Algebraic closure alone does not make
the resulting dynamics a gradient method in an arbitrarily assigned mirror.

In dual coordinates $y=\nabla h(x)$, write
$G(y)=\nabla h(F(\nabla h^*(y)))$. The equivalent symmetry test is
\begin{equation}
 DG(y)H_h(x)=H_h(x)DG(y)^T.
 \label{eq:dual-compatibility}
\end{equation}
For an affine dual map $G(y)=Ay+c$, this must hold for the entire family
of mirror Hessians, not just the Hessian at an anchor. For example,
$A=\diag(1/2,4/5)$ passes with $H_h=I$ and fails with
$H_h=\left[\begin{smallmatrix}2&0.6\\0.6&1\end{smallmatrix}\right]$.
Conversely, every Legendre geometry admits the family
$f=\alpha h-b^Tx$, with $G(y)=(1-\eta\alpha)y+\eta b$ whenever the image
stays in the dual domain. Thus arbitrary curvature can coexist with a
one-mode affine recurrence. More modes require a compatible Hessian family.

These statements refer to a prescribed mirror algorithm in its affine
coordinates. An arbitrary nonlinear coordinate change conjugates an existing
map, but recomputing ordinary mirror descent from a transformed objective
and transformed potential need not produce that conjugate map.

\subsection{The polynomial belongs to an observable recurrence}
Let $M$ be a state space and $F:M\to M$ any deterministic map. Define
$U\phi=\phi\circ F$. For a vector observable $\psi$, suppose
\begin{equation}
 \psi(F(z))=T\psi(z),\qquad \chi(\psi(z))=z
 \label{eq:general-lift}
\end{equation}
throughout an invariant set. Then
\begin{equation}
 F^m(z)=\chi(T^m\psi(z)).
 \label{eq:general-transport}
\end{equation}
If $\pi$ is the minimal polynomial of $T$ on the cyclic subspace generated
by $\psi(z)$, $T^m\psi(z)=p_m(T)\psi(z)$, where
$p_m(t)=t^m\bmod\pi(t)$. This is exact finite polynomial transport, including
unit eigenvalues and Jordan blocks, without a fixed-point inverse.
The decoder may be nonlinear. Requiring the original coordinate functions
to lie in the linear observable span would impose a stronger restriction.

For a coordinate change $\theta$, the observable
$\widetilde\psi=\psi\circ\theta^{-1}$ satisfies the same recurrence for
$\widetilde F=\theta F\theta^{-1}$, with decoder $\theta\circ\chi$.
Consequently existence of an injective finite lift is invariant under
conjugacy, but low degree in a chosen coordinate dictionary is not. Arnoldi
on a frozen tangent does not discover this property by itself.

A second route keeps the original dual coordinate as an additive output.
Set $g(y)=\nabla f(\nabla h^*(y))$. If the vector-valued functions
$g,Ug,U^2g,\ldots$ span a finite-dimensional space, then
\begin{equation}
 y_m=y_0-\eta\sum_{j=0}^{m-1}(U^jg)(y_0)
 \label{eq:gradient-module}
\end{equation}
can be evaluated through its discovered recurrence. This is a sufficient
condition for additive dual transport, not a necessary condition for the
nonlinear-decoder formulation \eqref{eq:general-lift}.

\subsection{A genuinely nonlinear dual map with finite closure}
For a scalar mirror let $s=h'(q)>0$ and choose $0<a\leq1$, $b\geq0$.
Define the objective by its derivative, with unit step size,
\begin{equation}
 f'(q)=s-\frac{as}{1+bs}.
\end{equation}
Its mirror update and Hessian are
\begin{equation}
 s^+=\frac{as}{1+bs},\qquad
 f''(q)=\left(1-\frac{a}{(1+bs)^2}\right)h''(q)\geq0.
 \label{eq:mobius-mirror}
\end{equation}
On an invariant positive dual interval this is convex mirror descent.
Despite nonlinear dual dynamics, $r=1/s$ obeys
\begin{equation}
 r^+=a^{-1}r+b/a.
 \label{eq:reciprocal-closure}
\end{equation}
The lift $(1,r)$ closes exactly and decoding is $q=(h')^{-1}(1/r)$.
The identity is obtained by division, so its global validity on $s>0$ is
structural; no dense trajectory sampling is required to prove it.

For $q=Cx$ with invertible $C$, let
$h(x)=\sum_i h_i(q_i)$ and $f(x)=\sum_i f_i(q_i)$.
The transformed dual coordinates are $s=C^{-T}\nabla h(x)$.
Equations \eqref{eq:mobius-mirror}--\eqref{eq:reciprocal-closure} apply componentwise.
This gives coupled primal geometries and arbitrary scalar mirror factors
whose domains support the chosen invariant interval. It is a separable
construction in known $q$ coordinates; it is not a theorem that every coupled
mirror permits independent prescribed rates. The compatibility theorem
specifies that remaining restriction.

If there are $p$ distinct values of $a_i$, the orbit of the augmented
reciprocal observable has rank at most $p+1$, independent of ambient dimension.
Rates need not be supplied to numerical recurrence identification. A dictionary
containing $s$, $\log s$, and $1/s$ can be screened by snapshot rank and
held-out defects. Providing that dictionary is prior representation knowledge;
finding it is distinct from fitting its recurrence.

We can go one step further within a finite rational grammar. From paired
observations $(s_j,t_j)$ fit a homogeneous relation
\begin{equation}
 [s_j,1,-s_jt_j,-t_j]\,(A,B,C,D)^T=0.
\end{equation}
When it identifies $t=(As+B)/(Cs+D)$, its fixed points are the roots of
$Cs^2+(D-A)s-B=0$. If these are distinct real finite roots $r_1,r_2$, define
\begin{equation}
 \phi(s)=\frac{s-r_1}{s-r_2},\qquad
 \phi(t)=\lambda\phi(s),\qquad
 \lambda=\frac{A-Cr_1}{A-Cr_2}.
 \label{eq:synthesized-crossratio}
\end{equation}
The identity follows by factoring $t-r_i$ and is valid away from poles.
The decoder is $s=(\phi r_2-r_1)/(\phi-1)$.
For the positive control family, choose the attracting fixed point $r_1=0$;
then $r_2=(a-1)/b<0$ and $\lambda=a$. We use $a<1,b>0$ in this experiment.
Thus the coordinate is synthesized from the relation, rather than supplied
as a reciprocal dictionary entry. The implementation rejects unidentified
relations and repeated, infinite, or nonreal roots; it does not claim to cover
those cases. A small fitted residual still needs an independent structural
certificate before being treated as a global identity.

\subsection{What discovery can certify, and what it cannot}
There are three separate claims: low rank of observed snapshots, a recurrence
valid on the future invariant set, and economical execution. SVD or Arnoldi
addresses the first. A symbolic identity such as
\eqref{eq:reciprocal-closure}, a certified residual bound on a reachable tube,
or unisolvent evaluation in a known finite residual space can address the
second. None alone establishes the third. A short observed orbit may miss
unexcited modes or future changes; numerical tolerances and conditioning
matter even when the underlying identity is exact.

Even Euclidean geometry does not guarantee finite linear closure of coordinate
observables. In one dimension, $h(x)=x^2/2$, $f(x)=x^4/4$, and $\eta=1$ give
$F(x)=x-x^3$. The polynomials $U^jx=F^j(x)$ have distinct degrees $3^j$ and
nonzero leading coefficients, hence are linearly independent on every open
interval. On $(-1/\sqrt3,1/\sqrt3)$ the map is invariant and monotone and the
objective is convex; the obstruction does not require erratic dynamics.
This rules out finite invariant linear spaces containing $x$ for this map.
It does not rule out a nonlinear decoder, a restricted orbit, or approximate
finite-horizon transport. Finite black-box samples also cannot certify
unrestricted smooth maps: a smooth perturbation supported away from all
queried states can alter later evolution while leaving every observed value
and derivative unchanged. Therefore a useful universal framework must expose
its structural assumptions rather than infer them from rank alone.

\subsection{Approximate transport without discarding directions}
For an imperfect lift define $\epsilon(z)=\psi(F(z))-T\psi(z)$.
On the actual trajectory $z_j=F^j(z_0)$, exact telescoping gives
\begin{equation}
 \psi(z_m)-T^m\psi(z_0)
   =\sum_{j=0}^{m-1}T^{m-1-j}\epsilon(z_j).
 \label{eq:directional-lift-defect}
\end{equation}
This retains directions and cancellation before any norm is taken. Replacing
it immediately by a sum of norms may destroy the structure that makes a
low-order recurrence useful. Its right side involves unknown actual states;
it is an identity, not a cheap certificate. Useful certification must enclose
those residuals over an invariant set or a justified reachable tube without
replaying the expensive evolution.

If both lifted endpoints lie in a region where $\chi$ is $L_\chi$-Lipschitz,
state error is at most $L_\chi$ times the norm of this vector sum. When the
transported coordinate is dual $y$ and its error is $e$, mirror duality gives
\begin{equation}
 D_h(x,\widehat x)=D_{h^*}(\widehat y,y)
 =\int_0^1(1-t)e^TH_{h^*}(y+te)e\,dt,
 \quad \widehat y=y+e.
 \label{eq:geometric-readout}
\end{equation}
This formula requires the dual segment to remain in the differentiability
domain. It separates recurrence error from geometric readout and does not
assume a global Euclidean Lipschitz bound. Near a singular decoder or boundary,
a small lifted residual may have a large state effect. Conversely, reciprocal
coordinates may have growing magnitude while decoding remains stable.

\subsection{Controlled cross-geometry discovery experiment}
We use 24-dimensional Euclidean, quartic, entropy, and hypentropy mirrors,
with a fixed dense invertible mixing of primal coordinates, three seeds,
and three distinct rates. The scalar mirror gradients are $q$, $q+q^3$,
$\log q$, and $\operatorname{arsinh}q$, respectively. Eight ordinary mirror
transitions supply nine snapshots; no future reference or rates enter fitting.
The known mirror factorization supplies the $s$ coordinates. The supplied
dictionary and synthesized rational chart are evaluated at horizons 16, 64,
and 128 against independent ordinary execution. Reported errors are maximum
relative primal errors across all 12 geometry/seed combinations.
\begin{center}\small
\begin{tabular}{lrrrr}\toprule
Representation & Snapshot rank & $m=16$ & $m=64$ & $m=128$\\\midrule
Dual coordinate & 8 & 1.08e-05 & 1.47e-01 & 4.53e+00\\
Log dual coordinate & 8 & 4.61e-06 & 4.34e-03 & 4.13e-02\\
Reciprocal dual coordinate & 4 & 1.71e-13 & 2.15e-10 & 2.45e-08\\
Synthesized cross-ratio & --- & 2.83e-14 & 2.78e-12 & 2.88e-11\\
\bottomrule\end{tabular}\end{center}
The SVD rank threshold is $10^{-10}$ relative to the largest singular value.
Cross-ratio synthesis fits a separate rational relation per transformed
coordinate and does not use a low-rank snapshot compression. Its rates are
learned, not supplied. The table measures a favorable constructed family;
it is not an acceptance-coverage study over arbitrary objectives. Numerical
errors reflect identification, decoding, and floating-point evolution, whereas
the exact closure statements follow from the algebra above. This experiment
records acquisition, fitting, transport, and ordinary execution time, but does
not perform a repeated performance gate; no wall-clock acceleration claim
is drawn from it. Ten new tests check mirror inversion, actual mirror updates,
convex Hessians, compatibility and its failure, recurrence recovery, chart
synthesis and rejection, directional defect cancellation, and polynomial
degree growth. The complete research suite passes 51 tests on the M4 Mini.


\subsection{The resulting research target}
The general object is a geometry-compatible map with an economical observable
recurrence and controlled decoding. A frozen tangent is one candidate
representation. The central task is to generate small observable families
from available mirror/proximal expressions, certify their residuals on useful
regions, and amortize discovery. Our control recovers a supplied reciprocal
chart and synthesizes a cross-ratio chart within a rational grammar; unrestricted
representation discovery remains open. The compatibility theorem and
geometry-sensitive error identity specify what broader constructions must prove.
